\begingroup
\setlength{\LTcapwidth}{\textwidth}
\begin{longtable}{p{0.34\textwidth}p{0.51\textwidth}}
\caption{Struktur informasi clinical state untuk integrasi antarmodul.}
\label{tab:tab5.7_struktur_clinical_state} \\
\toprule
\textbf{Elemen} & \textbf{Peran} \\
\midrule
\endfirsthead

\multicolumn{2}{@{}l}{\small\emph{Tabel \thetable{} (lanjutan)}} \\
\toprule
\textbf{Elemen} & \textbf{Peran} \\
\midrule
\endhead

\bottomrule
\endlastfoot

Sumber glukosa & Menentukan modalitas pengukuran yang digunakan \\

Nilai glukosa terakhir & Konteks keadaan terakhir yang diketahui \\

Prediksi glukosa & Nilai prediksi pada horizon tujuan \\

Horizon prediksi & Menentukan waktu yang direpresentasikan oleh prediksi \\

Interval ketidakpastian & Rentang prediksi apabila artefak kalibrasi tersedia \\

Kondisi klinis & Kategori kondisi hasil artefak klasifikasi yang sesuai \\

Konteks temporal & Informasi waktu dan perubahan yang diperlukan downstream \\
\end{longtable}
\endgroup
